\subsection{Kết quả đạt được}
Xây dựng thành công các mô hình dự đoán mức tiêu thụ năng lượng của các thiết bị trong nhà, với mô hình Random Forest đạt hiệu suất tốt nhất, cũng như xác định được yếu tố đóng vai trò quan trọng trong việc dự đoán. Song song với đó, bài báo cáo này là cơ hội tuyệt vời để phát triển được những kỹ năng phân tích, xử lý dữ liệu và học được cách sử dụng các thư viện của python để phân tích dữ liệu, xây dựng mô hình, đánh giá mô hình, vẽ biểu đồ,... Quan trọng hơn cả là áp dụng được các kiến thức đã học ở học phần Máy học ứng dụng vào bài toán thực tế.

Dữ liệu và code đã được công khai tại:
\url{https://github.com/wutdahack1337/CT294-project}

\subsection{Hướng phát triển}
Trong tương lai chúng ta có thể thử nghiệm thêm các mô hình học máy tiên tiến hơn như Deep Neural Network,... để nắm bắt được thêm các mối quan hệ phức tạp khác trong tập dữ liệu. Hoặc có thể thu thập thêm dữ liệu từ nhiều nguồn và thu thập thêm nhiều đặc trưng khác như lượng mưa, sự hiện diện của con người,... để xây dựng mô hình dự đoán toàn diện hơn.